\exercisesection{Local cohomology, Grothendieck duality}

\begin{enumerate}
\def\labelenumi{\arabic{enumi})}
\tightlist
\item
  Let A be a ring, J an ideal of A, M an A-module. We denote by \(H_{A,J}(M)\) , or simply \(H_{J}(M)\) , the graded A-module \(\varinjlim_{n}\operatorname{Ext}_{A}(A/J^{n},M)\) . If A is local Noetherian, the graded A-module \(H_{A,m_{A}}(M)\) is identified with \(H_{A}(M)\) .
\end{enumerate}

\begin{enumerate}
\def\labelenumi{\alph{enumi})}
\item
  Establish for \(H_{J}(M)\) the properties analogous to those of \(H_{A}(M)\) with respect to homomorphisms and exact sequences of A-modules (no. 1).
\item
  Define for every \(i \geqslant 1\) isomorphisms \(\mathrm{H}_{\mathrm{J}}^{i+1}(\mathrm{M}) \to \varinjlim \operatorname{Ext}_{\mathrm{A}}^{i}(\mathrm{J}^{n}, \mathrm{M})\) .
\item
  We suppose the ring A is local Noetherian of dimension d, and the ideal J is generated by a completely secant sequence \((x_{1},\ldots,x_{r})\) . Construct an isomorphism

\[
\tau^ {i} (\mathrm{J}; \mathrm{M}): ^ {\prime} \operatorname{Tor} _ {r - i} ^ {\mathrm{A}} (\mathrm{M}, \mathrm{H} _ {\mathrm{J}} ^ {d} (\mathrm{A})) \longrightarrow \mathrm{H} _ {\mathrm{J}} ^ {i} (\mathrm{M})
\]

which generalizes the isomorphism \(\tau^{i}(\mathrm{M})\) of no. 2; in particular, we have \(\mathrm{H}_{\mathrm{J}}^{i}(\mathrm{A}) = 0\) for \(i > r\) .

\item
  For every ideal K of A contained in J, we denote by \(\rho_{\mathrm{K},\mathrm{J}}:\mathrm{H}_{\mathrm{J}}(\mathrm{M})\to \mathrm{H}_{\mathrm{K}}(\mathrm{M})\) the homomorphism deduced from the canonical maps of \(\mathrm{A / K^n}\) in \(\mathrm{A / J^n}\) . Let I be an ideal of A; construct a long exact sequence of A-modules

\[
\dots \longrightarrow \mathrm{H} _ {\mathrm{I} +, \mathrm{J}} ^ {i} (\mathrm{M}) \stackrel {\alpha} {\longrightarrow} \mathrm{H} _ {\mathrm{I}} ^ {i} (\mathrm{M}) \oplus \mathrm{H} _ {\mathrm{J}} ^ {i} (\mathrm{M}) \stackrel {\beta} {\longrightarrow} \mathrm{H} _ {\mathrm{I} \cap \mathrm{J}} ^ {i} (\mathrm{M}) \longrightarrow \mathrm{H} _ {\mathrm{I} \mid \mathrm{J}} ^ {i + 1} (\mathrm{M}) \longrightarrow \dots ,
\]

\[
\text {   où   } \alpha (x) = (\rho_ {\mathrm{I,I+J}} (x), \rho_ {\mathrm{J,I+J}} (x)) \text {   et   } \beta (x, y) = \rho_ {\mathrm{I} \cap \mathrm{J}, \mathrm{I}} (x) - \rho_ {\mathrm{I} \cap \mathrm{J}, \mathrm{J}} (y).
\]

\end{enumerate}

\begin{enumerate}
\def\labelenumi{\arabic{enumi})}
\setcounter{enumi}{1}
\tightlist
\item
  Let \(\Lambda\) a local Macaulay ring of dimension \(d\) . Prove that for A to be a Gorenstein ring, it is necessary and sufficient that the \(\Lambda\) -module \(\mathrm{H}_{\mathrm{A}}^{d}(\Lambda)\) be injective.
\end{enumerate}

¶ 3) Let \(\rho: A \to B\) be a local homomorphism of local Noetherian rings, such that \(\rho(\mathfrak{m}_{A})B\) is a definition ideal of B.

\begin{enumerate}
\def\labelenumi{\alph{enumi})}
\item
  Construct for every B-module N a natural A-linear isomorphism of \(H_{A}(N)\) on \(H_{B}(N)\) . (Using an injective resolution of N, show that it suffices to prove that we have \(H_{A}^{i}(J)=0\) for all i\textgreater0 and every indecomposable injective B-module J. Let q be the unique element of \(\operatorname{Ass}_{\mathrm{B}}(\mathrm{J})\) , and by \(p=\rho^{-1}(q)\) If \(p\neq m_{A}\) , observe that there exists \(a\in m_{A}\) such that the homothety \(a_{J}\) be bijective. If \(p=m_{A}\) , and if I is a minimal injective resolution of the A-module J, observe that we have \(\operatorname{Ass}_{\mathrm{A}}(\mathrm{I}^{n})=\{\mathfrak{m}_{\mathrm{A}}\}\) for every n.)
\item
  We suppose moreover that B is a flat A-module. Construct for every A-module M a natural B-linear isomorphism of \(\mathrm{B}\otimes_{\mathrm{A}}\mathrm{H}_{\mathrm{A}}(\mathrm{M})\) on \(\mathrm{H}_{\mathrm{B}}(\mathrm{B}\otimes_{\mathrm{A}}\mathrm{M})\) .
\end{enumerate}

\begin{enumerate}
\def\labelenumi{\arabic{enumi})}
\setcounter{enumi}{3}
\item
  Let A be a local ring and M a nonzero finitely generated A-module, of dimension e. Prove that \(\mathrm{H}_{\mathrm{A}}^{e}(\mathrm{M})\) is not zero (reduce to the case where the ring A is complete, then with the aid of exerc. 3 to the case where A is regular).
\item
  Let I be a finite subset of N.
\end{enumerate}

\begin{enumerate}
\def\labelenumi{\alph{enumi})}
\item
  Construct a regular local ring B and a B-module N such that the set of centers p such that \(H_{\mathrm{B}}^{p}(\mathrm{N}) \neq 0\) is equal to I (consider a direct sum).
\item
  Construct a local ring A such that the set of integers p such that \(\mathrm{H}_{\mathrm{A}}^{p}(\mathrm{A}) \neq 0\) is equal to I (take \(A = B \oplus N\) , where N is an ideal of square zero, and use exerc. 3).
\end{enumerate}

\begin{enumerate}
\def\labelenumi{\arabic{enumi})}
\setcounter{enumi}{5}
\tightlist
\item
  Let A be a local Noetherian ring admitting a dualizing module, M a finitely generated A-module, n an integer.
\end{enumerate}

\begin{enumerate}
\def\labelenumi{\alph{enumi})}
\item
  For the A-module \(H_{\mathrm{A}}^{n}(\mathrm{M})\) to be of finite length, it is necessary and sufficient that for every prime ideal p of A distinct from \(m_{A}\) , one has \(\mathrm{H}_{\mathrm{A}_{\mathfrak{p}}}^{n-\dim(\mathrm{A}/\mathfrak{p})}(\mathrm{M}_{\mathfrak{p}})=0\) (use Grothendieck duality).
\item
  For \(\mathrm{H}_{\mathrm{A}}^{i}(\mathrm{M})\) be of finite length for every \(i \leqslant n\) , it is necessary and sufficient that one has \(\operatorname{prof}(\mathrm{M}_{\mathfrak{p}}) + \dim(\mathrm{A}/\mathfrak{p}) > n\) for every prime ideal p distinct from \(m_{A}\)
\end{enumerate}

\begin{enumerate}
\def\labelenumi{\arabic{enumi})}
\setcounter{enumi}{6}
\tightlist
\item
  Let A be a ring, M a \(\Lambda\) -module, \(\mathbf{x} = (x_i)_{i \in I}\) a finite family of elements of A. For every integer \(n \geqslant 1\) , we denote \(\mathbf{x}^n\) the family \((x_i^n)_{i \in I}\) . Let \(\Delta(\mathbf{x})\) the endomorphism of \(\mathrm{A}^{\mathrm{I}}\) defined by \(\Delta(\mathbf{x})(e_i) = x_i e_i\) .
\end{enumerate}

\begin{enumerate}
\def\labelenumi{\alph{enumi})}
\tightlist
\item
  For every pair of integers \((n, m)\) such that \(n \geqslant m\) show that the maps
\end{enumerate}

\[
\operatorname{Homgr} \left(\Lambda \left(\Delta \left(\mathbf {x} ^ {n - m}\right)\right), 1 _ {\mathrm{M}}\right): \mathbf {K} ^ {\bullet} \left(\mathbf {x} ^ {m}, \mathrm{M}\right) \longrightarrow \mathbf {K} ^ {\bullet} \left(\mathbf {x} ^ {n}, \mathrm{M}\right)
\]

define an inductive system of complexes; we denote by \(\mathbf{K}^{\bullet}(\mathbf{x}^{\infty},\mathrm{M})\) the limit of this system, and \(\mathrm{H}^{\bullet}(\mathbf{x}^{\infty},\mathrm{M})\) its cohomology.

\begin{enumerate}
\def\labelenumi{\alph{enumi})}
\setcounter{enumi}{1}
\tightlist
\item
  We choose a total order on I. For every subset J of I, we set \(x_{J} = \prod_{j \in J} x_{j}\) .
\end{enumerate}

For \(p \in \mathbf{Z}\) , we set \(\mathcal{C}^p = \bigoplus_{|J| = p} A_{x_J}\) ; for \(a \in A_{x_J}\) , with \(\text{Card}(J) = p\) , we set \(d^p(a) = \sum_{i \notin J} (-1)^{\varepsilon(J, i)} \frac{a}{1}\) in \(A_{x_{J \cup \{i\}}}\) , where \(\varepsilon(J, i)\) is the number of elements of \(J\) strictly less than \(i\) Prove that \((\mathcal{C}, d)\) is a complex, and define a canonical isomorphism of \(\mathbf{K}^\bullet(\mathbf{x}^\infty, M)\) on the complex \(\mathcal{C} \otimes_A M\) .

8) We keep the notation from the previous exercise, assuming moreover that the ring A is Noetherian. We denote by x the ideal generated by x.

\begin{enumerate}
\def\labelenumi{\alph{enumi})}
\item
  Let m be an integer. Prove that there exists an integer \(n_{0} \geqslant m\) such that for \(n \geqslant n_{0}\) such that the morphism \(\boldsymbol{\Lambda}(\Delta(\mathbf{x})^{n-m}) : \mathbf{K}_{\bullet}(\mathbf{x}^{n}, \mathrm{A}) \longrightarrow \mathbf{K}_{\bullet}(\mathbf{x}^{m}, \mathrm{A})\) induces 0 on the homology in degree \(\geqslant 1\) (observe that one has \(\boldsymbol{\Lambda}^{p}(\Delta(\mathbf{x}^{n-m}))(\mathbf{K}_{p}(\mathbf{x}^{n}, \mathrm{A})) \subset x^{n-m}\mathbf{K}_{p}(\mathbf{x}^{m}, \mathrm{A})\) ; deduce from the Artin–Rees lemma that there exists an integer \(n_{0}\) such that \(Z_{p}(\mathbf{x}^{m}, \mathrm{A}) \cap x^{n_{0}}\mathbf{K}_{\bullet}(\mathbf{x}^{m}, \mathrm{A})\) is contained in \(x^{rm}Z_{p}(\mathbf{x}^{m}, \mathrm{A})\) , and hence in \(B_{p}(\mathbf{x}^{m}, \mathrm{A})\) ).
\item
  Define a canonical isomorphism from \(\varinjlim_{n}\mathrm{Ext}_{\mathrm{A}}^{p}(\mathrm{A}/x^{n},\mathrm{M})\) on \(\mathrm{H}^p (\mathbf{x}^\infty ,\mathrm{M})\) (reduce to proving that \(\mathrm{H}^p (\mathbf{x}^\infty ,\mathrm{M})\) is zero for \(p\geqslant 1\) and M injective; in this case observe that \(\mathrm{H}^p (\mathbf{x}^\infty ,\mathrm{M})\) is identified with the inductive limit of the \(\mathrm{Hom}_{\mathrm{A}}(\mathrm{H}_p(\mathbf{x}^n,\mathrm{M}))\) and use \(a)\) ). In particular, if A is local and if \(x\) is an ideal of definition of A, the graded A-module \(\mathrm{H}_{\mathrm{A}}(\mathrm{M})\) is canonically isomorphic to \(\mathrm{H}^{\bullet}(\mathbf{x}^{\infty},\mathrm{M})\) .
\item
  We assume the ring A is local Noetherian, of dimension \(d\) ; let \(\mathbf{x} = (x_{1},\ldots ,x_{d})\) a maximal secant sequence of elements of \(\mathfrak{m}_{\Lambda}\) Show that the A-module \(\mathrm{H}_{\mathrm{A}}^{d}(\mathrm{M})\) is identified with the limit of the inductive system of A-modules \((\mathrm{M}_p,u_{qp})\) , where \(\mathrm{M}_p\) is the quotient module \(\mathrm{M} / (x_1^p\mathrm{M} + \dots +x_d^p\mathrm{M})\) and \(u_{qp}:\mathrm{M}_p\to \mathrm{M}_q\) ( \(q\geqslant p\) ) the homomorphism induced by the homothety of ratio \((x_{1}\dots x_{d})^{q - p}\) .
\end{enumerate}

\begin{enumerate}
\def\labelenumi{\arabic{enumi})}
\setcounter{enumi}{8}
\tightlist
\item
  Let A be a Noetherian local ring, d its dimension, \(\mathbf{x} = (x_{1}, \ldots, x_{d})\) a maximal secant sequence of elements of \(m_{A}\) .
\end{enumerate}

\begin{enumerate}
\def\labelenumi{\alph{enumi})}
\item
  Prove that there exists an integer \(m_0\) such that for \(m \geqslant m_0\) and \(n \geqslant 0\) , one has \((x_1^m \ldots x_d^m)^n \not\in (x_1^{m(n+1)}, \ldots, x_d^{m(n+1)})\) (use Exercises 4 and 7, b)).
\item
  Assume that \(\Lambda\) contains a field of characteristic \(p\) . Prove that we have \((x_{1} \ldots x_{d})^{n} \notin (x_{1}^{n+1}, \ldots, x_{d}^{n+1})\) for every \(n \geqslant 0\) , that is, A satisfies condition (CM) (Exercise 3 of § 2; consider the endomorphism \(x \mapsto x^p\) Therefore, every Noetherian local ring containing a field satisfies condition (CM). \(^{1}\) .
\item
  We further assume that A is regular; let \(\rho: A \to B\) an injective ring homomorphism making B into a finitely generated A-module. Prove that the A-module \(\rho(A)\) is a direct factor in B (cf.~§ 2, exerc. 3); reduce to the case where B is integral, hence local, and apply exerc. 8 of § 4).
\end{enumerate}

\begin{enumerate}
\def\labelenumi{\arabic{enumi})}
\setcounter{enumi}{9}
\tightlist
\item
  Let \(k\) a field, A a \(k\) -regular local algebra of dimension \(d\) , such that the extension \(k \to \kappa_{\mathrm{A}}\) is trivial. The A-modules \(\mathrm{H}_{\mathrm{A}}^{d}(\mathrm{A})\) and \(\mathrm{A}' = \mathrm{Hom}_k^{cont}(\mathrm{A}, k)\) are Matlis modules, hence isomorphic; we propose to make explicit an isomorphism between these two modules.
\end{enumerate}

\begin{enumerate}
\def\labelenumi{\alph{enumi})}
\item
  Let \(\mathbf{x} = (x_{1}, \ldots, x_{d})\) a coordinate system of A. For every integer \(p \geqslant 0\) , we denote \(A_{p}\) the k-algebra \(\mathrm{A}/(x_{1}^{p}, \ldots, x_{d}^{p})\) The classes of monomials \(x_{1}^{a_{1}} \ldots x_{d}^{a_{d}}\) with \(0 \leqslant a_{i} < p\) form a basis of \(A_{p}\) on k; we denote by \(\rho_{p}\) the linear form on \(A_{p}\) which takes the value 1 on the class of \(x_{1}^{p-1} \ldots x_{d}^{p-1}\) and 0 on the other elements of the basis. Prove that the k-linear map \(\tilde{\rho}_{p}: A_{p} \to \operatorname{Hom}_{k}(A_{p}, k)\) derived from the bilinear form \((a, b) \mapsto \rho_{p}(ab)\) is an isomorphism (one may reason directly, or use Exercise 10 of § 5).
\item
  For \(q \geqslant p\) let \(u_{qp}: A_p \to A_q\) the A-linear homomorphism induced by the homothety of ratio \((x_1 \ldots x_d)^{q-p}\) , and let \(\pi_{pq}: A_q \to A_p\) the passage-to-the-quotient map; prove that one has \(\tilde{\rho}_q \circ u_{qp} = {}^t \pi_{pq} \circ \tilde{\rho}_p\) . Deduce by passage to the inductive limit, taking into account exerc. 8, c), an A-linear isomorphism \(\tilde{\rho}: H_A^d(A) \to A'\) .
\item
  Assume \(d=1\) , so that A is a discrete valuation ring, and \(x_{1}=x\) a uniformizer of A. Show that the choice of the uniformizer \(x\) allows one to identify \(\mathrm{H}_{\mathrm{A}}^{1}(\mathrm{A})\) to K/A; the isomorphism \(\rho:\mathrm{K}/\mathrm{A}\to\mathrm{A}'\) is then defined by \(\rho(\varphi)(a)=\mathrm{Res}(a\varphi)\) , where \(\mathrm{Res}:\mathrm{K}/\mathrm{A}\to k\) is the projection onto the factor \(k x^{-1}\) in the canonical decomposition \(\mathrm{K}/\Lambda=\bigoplus_{n\geqslant 1}k x^{-n}\) .
\end{enumerate}

\begin{enumerate}
\def\labelenumi{\arabic{enumi})}
\setcounter{enumi}{10}
\item
  Let A be a Noetherian local ring, I a Matlis module over A, and M a finitely generated A-module. We set \(\mu^{p} = \dim_{\kappa_{A}} \operatorname{Ext}_{A}^{p}(\kappa_{A}, M)\) for every \(p \geqslant 0\) (cf. § 8, exerc. 4). Prove that the graded A-module H \(_{A}\) (M) is identified with the homology of a complex \(\ldots \to 0 \to I^{\mu^{0}} \to I^{\mu^{1}} \to \ldots\) (loc. cit.). In particular, if M is nonzero, we have \(\mu^{p} \neq 0\) for \(p = \text{prof}(M)\) and \(p = \dim(M)\) (exercise 4).
\item
  Let A be a Noetherian local ring, M a finitely generated A-module of dimension d; suppose that M is a Buchsbaum module (§ 2, exerc. 7).
\end{enumerate}

\begin{enumerate}
\def\labelenumi{\alph{enumi})}
\item
  Prove that \(\mathrm{H}_{\mathrm{A}}^{i}(\mathrm{M})\) is annihilated by \(m_{A}\) for \(i \neq d\) (First handle the case d = 1, then reason by induction on d).
\item
  Let \(x\) a secant element for \(\mathfrak{m}\) ; define for \(0 \leqslant i \leqslant d - 2\) the exact sequences \(0 \to \mathrm{H}_{\mathrm{A}}^{i}(\mathrm{M}) \to \mathrm{H}_{\mathrm{A}}^{i}(\mathrm{M}/x\mathrm{M}) \to \mathrm{H}_{\mathrm{A}}^{i+1}(\mathrm{M}) \to 0\) .
\item
  Prove the formula \(i(\mathrm{M}) = \sum_{i=0}^{d-1} \left( \begin{array}{cc} d & 1 \\ i & \end{array} \right) \lg(\mathrm{H}_{\mathrm{A}}^i(\mathrm{M}))\) (argue by induction on the integer \(d \geqslant 1\) , using \(b\) ) and exerc. 8, \(d\) ) of § 2).
\end{enumerate}

